% 日常任务的完整证据链 -> 手术模拟任务上的跨任务适用性。
% 所有结果均留空；括号内下划线文字是待填写的分析提示。
% 组织参考：Cable Routing (T-RO 2024) 的研究问题与分层验证，GC-DA (T-RO 2024) 的问题驱动消融。
\section{Experimental Evaluation}
\label{sec:experiments}
This section evaluates the proposed method on a bimanual exposure and execution task.
The task requires coordinated actions under deformable occlusion,
where preceding interactions change the visibility and accessibility of the target.
We first assess overall task performance across object categories, initial positions, and occluder colors,
then examine annotation reliability, semantic representations, and history-conditioned correction through controlled ablations.
Finally, we evaluate the framework on a surgical-phantom task to assess its applicability beyond the primary setting.

主要实验围绕四个问题展开：
\begin{enumerate}
\item 完整方法能否提高任务完成率，且收益是否在不同物体、初始位置和遮挡物颜色上保持一致？
\item 质量感知迭代能否在保持分割精度的同时减少异常，并降低人工标注负担？
\item 语义输入与任务描述量监督分别带来什么收益，感知质量如何影响这些收益？
\item 更新观测与历史能否减少执行错误，并促进自然偏差后的恢复？
\end{enumerate}

\subsection{Experimental Setup}
\label{sec:evalprotocol}
% 主实验任务；附加手术模拟任务在跨任务适用性小节单独定义。
\subsubsection{Task and Platform}
平台采用两条SO-ARM101机械臂，右臂末端安装长板工具，
通过固定的前视与侧视相机观察桌面操作区域。
示范中的目标包括方块、纸团和螺丝螺母三类，遮挡物为同一颜色的可形变布料。
任务开始时目标位于布料下方，双臂协同拨开布料，使目标暴露并与布料分离，
再将目标推入指定区域，最后恢复布料在前视画面中的覆盖率。

以完整目标投影 $O_t^\star$、完整目标区域足迹 $R^\star$ 和布料投影 $C_t$
定义入区与复原条件：
\begin{equation}
 \begin{aligned}
 p_{\mathrm{in}}(t)&=\ind\{O_t^\star\subseteq R^\star\},\\
 p_{\mathrm{restore}}(t)&=
 \ind\{|C_t\cap\Omega|/|\Omega|\ge\theta_C\},
 \end{aligned}
 \label{eq:taskpredicates}
\end{equation}
其中 $\Omega$ 为固定前视相机的整幅图像域，$\theta_C$ 为布料覆盖阈值。
任务成功要求目标完整入区且布料覆盖率达到阈值，
两项条件须在规定时限内同时成立：
\begin{equation}
 \Phi(s_{0:T})=
 \ind\{\exists\,t\le T:
 p_{\mathrm{in}}(t)=1\land p_{\mathrm{restore}}(t)=1\}.
 \label{eq:completiontask}
\end{equation}
\draftnote{填写最大执行时间与布料覆盖阈值。}


\begin{figure*}[t]
\centering
\fbox{\parbox[c][2.1cm][c]{0.95\textwidth}{\centering
(a) 双臂平台全景、前视与侧视相机、工具及目标区域\\[4pt]
(b) 同一试验的初始遮挡--暴露--分离--入区--布料复原关键帧\\[4pt]
替换为真实照片；图中标明前序交互如何改变后续操作条件}}
\caption{Bimanual exposure and execution: the experimental platform and a complete task sequence.}
\label{fig:tasksetup}
\end{figure*}

\subsubsection{Demonstrations and Test Conditions}
数据集包含90条示范、53,723个帧对齐样本，采集频率为30\,Hz。
前视与侧视原始图像均为$640\times360$，本体状态和动作均为10维关节位置，
对应每条机械臂的5个关节。

物体初始位置按$3\times4$网格布置，示范覆盖其中六个位置，如图~\ref{fig:objectpositions}所示；
基础测试覆盖三类已见物体与全部12个位置的组合，共36个条件。
另使用未见颜色的遮挡布料，在相同物体与位置条件下检验外观变化的影响。
各方法在每个条件下进行相同次数的独立重置试验，分别报告已见与未见位置的表现。
\draftnote{确定未见布料颜色与重复次数；固定三类物体、12个位置和各布料颜色的测试组合。}

\begin{figure}[t]
\centering
\fbox{\parbox[c][2.8cm][c]{.92\columnwidth}{\centering
(a) 方块、纸团、螺丝螺母及已见／未见颜色布料\\[4pt]
(b) 桌面$3\times4$网格及固定目标区域\\[4pt]
实心：6个示范位置；空心：6个未见测试位置\\[4pt]
标注行列方向、间距和前视相机方向}}
\caption{Object categories, occluder colors, and initial positions. Filled and open markers indicate demonstration and unseen test positions, respectively.}
\label{fig:objectpositions}
\end{figure}

\subsubsection{Training Settings}
各方法在固定示范预算下使用相同的RGB观测、本体状态、动作接口和示范划分；语义方法增加由示范视频获得的语义监督。
基础策略匹配训练更新次数与样本访问量，语义预训练和残差训练的额外资源单独报告。
残差消融共用冻结的基础策略，模型选择与最终测试分开。
表~\ref{tab:protocol}汇总关键执行设置，数据使用与评价细则见附录~\ref{app:datascope}及\ref{app:evaluationdetails}。

\begin{table}[t]
\caption{Key training and execution settings.}
\label{tab:protocol}
\centering\small
\begin{tabular}{p{.38\columnwidth}p{.49\columnwidth}}
\toprule
项目 & 配置\\
\midrule
数据划分／训练种子 & 待填\\
每个条件执行重复 & 待填\\
时限／布料覆盖阈值 & 待填\\
控制频率／部署设备 & 待填\\
动作块／重规划间隔 & 待填\\
残差间隔／预留延迟 & 待填\\
\bottomrule
\end{tabular}
\end{table}

\subsubsection{Evaluation Metrics}
主要指标为固定时限内的完整任务成功率，报告成功次数／总次数及不确定性区间。
补充指标包括成功试验的完成时间、以全部试验为分母的累计成功时间曲线、
阶段到达与条件完成率，以及主要失败类型。
各方法采用一致的初始条件和计时规则，测试顺序在条件区组内随机安排。
任务成功与阶段进展由独立视频审核判定。

\subsection{Overall Task Performance}
\label{sec:overallperformance}
本节比较整体任务表现，并考察物体类别、初始位置和遮挡物颜色变化下的成功率。
对照包括ACT~\cite{zhao2023act}、Diffusion Policy~\cite{chi2023diffusion}和完整方法；
另预留A2C2~\cite{a2c2}作为在线动作块纠正对照、VIOLA~\cite{zhu2023viola}作为对象中心表示对照。
\draftnote{完成所选基线的合理调试后再纳入最终表格，报告适配改动与额外资源。
ACT上的A2C2适配不等于原版复现，本文MLP消融也不能直接标为A2C2。}

各基线采用匹配的模型选择与调参预算。
推理速度以部署设备上的端到端测量为准，不以名义更新频率代替实际频率。

\subsubsection{Task Success and Completion Time}
表~\ref{tab:mainresults}汇总三个已见物体类别、已见与未见位置、已见与未见遮挡物颜色的成功率，
以及整体成功率和成功试验的平均完成时间。
整体成功率以预先确定的完整测试池为分母，每次试验只计一次；分组列分别统计满足对应条件的试验。
图~\ref{fig:positionperformance}进一步展示逐位置表现。

\begin{table*}[t]
\caption{Task success across object categories, initial positions, and occluder colors, with overall success and mean completion time. Success entries report count/total and percentage; time is computed over successful trials.}
\label{tab:mainresults}
\centering\small
\setlength{\tabcolsep}{3pt}
\begin{tabular}{lccccccccc}
\toprule
& \multicolumn{3}{c}{类别成功率} & \multicolumn{2}{c}{位置成功率} & \multicolumn{2}{c}{遮挡物颜色成功率} & \multicolumn{2}{c}{整体表现}\\
\cmidrule(lr){2-4}\cmidrule(lr){5-6}\cmidrule(lr){7-8}\cmidrule(lr){9-10}
方法 & 方块 & 纸团 & 螺丝螺母 & 已见 & 未见 & 已见 & 未见 & 成功率 & 时间(s)\\
& 已见 & 已见 & 已见 & 6位置 & 6位置 & 颜色 & 颜色 & & \\
\midrule
ACT & -- & -- & -- & -- & -- & -- & -- & -- & --\\
Diffusion Policy & -- & -- & -- & -- & -- & -- & -- & -- & --\\
ACT＋A2C2适配 & -- & -- & -- & -- & -- & -- & -- & -- & --\\
VIOLA适配 & -- & -- & -- & -- & -- & -- & -- & -- & --\\
Ours & -- & -- & -- & -- & -- & -- & -- & -- & --\\
\bottomrule
\end{tabular}
\end{table*}

\begin{figure*}[t]
\centering
\fbox{\parbox[c][2.1cm][c]{.95\textwidth}{\centering
行：方块、纸团、螺丝螺母；列：ACT、DP、Ours\\[4pt]
每个面板为$3\times4$位置成功率热图，统一色标，标出示范位置\\[4pt]
每格标记成功次数／重复次数，按布料颜色分面；未测试组合留空}}
\caption{Position-wise task success by object category and occluder color. Highlighted cells indicate demonstration positions; untested combinations remain blank.}
\label{fig:positionperformance}
\end{figure*}
\draftnote{先比较整体成功率与完成时间，再检查收益在三个物体类别上是否一致，
以及未见位置和颜色下是否保持。结合全部试验的成功时间曲线，避免仅由成功样本均值判断效率。
分别分析位置、颜色变化及两者同时变化的情况；结论限定为固定示范预算下的表现。}

\subsubsection{Completion Time, Stage Progress, and Failures}
图~\ref{fig:timeandstages}比较完成时间与阶段到达情况，
表~\ref{tab:failurecounts}汇总主要失败类型。
二者用于区分暴露、目标执行与复原环节的改善，具体失败机制在Discussion中分析。
\begin{figure*}[t]
\centering
\fbox{\parbox[c][2.1cm][c]{0.95\textwidth}{\centering
(a) 截至时间$t$的完整成功比例曲线\\[4pt]
(b,c) 暴露、分离、入区、复原的累计到达率／条件完成率\\[4pt]
按方法绘制，包含误差区间和分母；不将失败试验从时间曲线中删除}}
\caption{Completion-time and stage-level comparisons. Placeholder for measured curves and stage statistics.}
\label{fig:timeandstages}
\end{figure*}

\begin{table}[t]
\caption{Primary failure categories for the everyday task. Report counts and fractions of all trials for each method.}
\label{tab:failurecounts}
\centering\small
\begin{tabular}{p{.22\columnwidth}p{.34\columnwidth}p{.29\columnwidth}}
\toprule
方法 & 主要失败类型 & 次数／全部试验\\
\midrule
待填 & 未完成暴露 & --\\
待填 & 无效接触／拨空气 & --\\
待填 & 目标移出可操作范围 & --\\
待填 & 入区不完整 & --\\
待填 & 复原不足 & --\\
待填 & 其他／原因未确定 & --\\
\bottomrule
\end{tabular}
\end{table}

\draftnote{比较完整方法相对各基线的成功率差值及区间，再结合时间曲线和阶段统计说明收益发生在哪个环节。
若只在某类物体或初始条件下改善，直接限定结论范围。表~\ref{tab:failurecounts}只统计现象，
具体原因结合视频与日志在Discussion分析；未完成的可选基线行在定稿时删除。}
\FloatBarrier

\subsection{Quality-Aware Semantic Annotation}
\label{sec:annotationexperiments}
本节检验质量感知迭代能否在保持分割精度的同时减少异常，并降低获得可靠标签的人工成本。
按完整视频划分标注开发数据与独立人工审核数据。
比较种子模型、无质量筛选的迭代和质量筛选迭代；
两条迭代路径使用相同候选视频池及匹配训练预算，并记录实际保留的数据量。
独立人工审核视频包含普通帧以及暴露、接触和入区等关键事件，
两类样本分层报告，评价不以迭代自身的质量分数代替人工参考。

表~\ref{tab:annotationquality}报告原始输出的mIoU、异常率、可用标签比例与人工投入。
异常帧由预先定义的漏分、错分和非真实结构突变规则审核；
时间一致性使用有效光流对应区域的运动补偿IoU。
原始输出与后处理结果分开统计，以区分模型本身的改善和筛选／修补的作用。
按类别补充目标和工具的IoU及异常计数，跨视频指标以视频为单位汇总。
人工成本计入关键帧标注、检查与修正时间，并结合达到统一审核标准的可用标签数量报告。
在相同视频子集上记录人工逐帧标注／修正流程的耗时，作为人工成本参照。
\begin{table*}[t]
\caption{Annotation quality and human effort on independently reviewed everyday-task videos. Iteration index $k$ is not fixed in advance.}
\label{tab:annotationquality}
\centering\small
\begin{tabular}{p{.22\textwidth}cccccc}
\toprule
标注／分割版本 & 标注视频数 & raw mIoU & 原始异常率 & 运动IoU & 保留标签率 & 人工分钟／视频\\
\midrule
种子模型 & -- & -- & -- & -- & -- & --\\
无质量筛选迭代 & -- & -- & -- & -- & -- & --\\
质量筛选迭代1 & -- & -- & -- & -- & -- & --\\
质量筛选迭代$k$ & -- & -- & -- & -- & -- & --\\
\bottomrule
\end{tabular}
\end{table*}

\begin{figure}[t]
\centering
\fbox{\parbox[c][3cm][c]{.92\columnwidth}{\centering
(a) 累积标注量--精度／异常率曲线\\[4pt]
(b) 同一视频的质量轨迹与真实事件标记\\[4pt]
(c) 异常帧的原始预测、处理结果和人工参考\\[4pt]
保留真实暴露／接触变化作为对照}}
\caption{Annotation expansion and temporal quality diagnostics. Placeholder for measured curves and reviewed mask examples.}
\label{fig:annotationquality}
\end{figure}
\draftnote{先比较无筛选与有筛选迭代，再分析每轮精度、异常率、保留率及人工成本的变化。
不能把排除了更多帧等同于模型更稳定；分别说明原始预测与后处理的贡献。
如精度基本持平而异常下降，应以稳定性改善为结论，不写成空间精度显著提升。
填写人工耗时记录与对照子集；没有实际耗时时仅报告异常和复核需求，暂不声称节省人工时间。}
\draftnote{核对用于该实验的分割训练记录。现有数据集的前视标注继承旧版本，侧视来自另一标注流程，
不能将已有标签直接记为最新迭代方法的产物；必要时在相同原始视频上补做匹配实验。
分别报告前视与侧视结果，人工成本同时计入种子标注和后续复核。}
\FloatBarrier

\subsection{Semantic Representation Ablations}
\label{sec:representationexperiments}
本节在不启用残差反馈的条件下，区分语义输入、查询容量和描述量监督的贡献。
各组在三类物体和12个位置上采用相同的测试条件，主要消融沿用已见颜色的布料。
正文报告分类别、已见／未见位置成功率以及阶段完成率和失败类型，完整逐条件结果放补充材料。
表~\ref{tab:representationablation}中S0--S3逐步加入语义与查询；
S2和S3保持查询数量及结构一致，只改变描述量辅助监督。
S3与S4使用相同标签和分割器，仅改变质量权重。
S4停用质量筛选与降权，但保留相同的几何有效性标记，不对无定义距离进行监督。
这些配置分别训练基础策略，并保持动作监督、视角及预算一致。

S5固定S3的策略权重，仅在推理时替换为验证集上预先选定的较低质量分割器，
保持类别定义、输入尺寸与语义编码一致，以评价在线感知质量的影响。
分割器更换造成的延迟变化单独记录。
此组不同时改变训练标签；不同质量标签训练出的策略属于另一种对照，不能混合解释。

\begin{table*}[t]
\caption{Representation ablations without residual feedback. UQ: queries without descriptor supervision; SQ: supervised task-descriptor queries.}
\label{tab:representationablation}
\centering\small
\begin{tabular}{clccccc}
\toprule
ID & 输入／查询 & 质量权重 & 推理分割 & 策略重训 & 成功率 & 端到端p95(ms)\\
\midrule
S0 & RGB & 不适用 & 无 & 是 & -- & --\\
S1 & RGB＋语义 & 不适用 & 标准 & 是 & -- & --\\
S2 & RGB＋语义＋UQ & 不适用 & 标准 & 是 & -- & --\\
S3 & RGB＋语义＋SQ & 启用 & 标准 & 是 & -- & --\\
S4 & RGB＋语义＋SQ & 停用 & 标准 & 是 & -- & --\\
S5 & 同S3 & 同S3 & 较低质量 & 否 & -- & --\\
\bottomrule
\end{tabular}
\end{table*}

图~\ref{fig:representationdiagnostics}按物体与位置汇总任务结果，并定位阶段收益与失败类型的变化。
进一步在独立审核数据上评价伪标签误差和部署模式下的描述量读出误差，
按面积、距离和方向分别报告，并保留有效标签比例。
补充的专家动作重构探针用于考察描述量是否保留动作选择所需的信息，
其划分与对照见附录~\ref{app:evaluationdiagnostics}。
该分析检验第~\ref{sec:representationanalysis}节提出的任务信息保留路径。

\begin{figure*}[t]
\centering
\fbox{\parbox[c][2.7cm][c]{.95\textwidth}{\centering
(a) S0--S5三个类别及已见／未见位置的成功率\\[4pt]
(b) 阶段累计到达率与条件完成率；(c) 主要失败类型占全部试验的比例\\[4pt]
(d) 伪标签与模型读出的分类描述量误差及有效标签比例\\[4pt]
标明样本量和区间；专家动作重构探针结果放附录}}
\caption{Representation ablations: task success, stage progress, failure categories, and descriptor accuracy.}
\label{fig:representationdiagnostics}
\end{figure*}

\draftnote{S1--S0检验语义输入；S3--S2检验描述量监督是否超出查询容量本身的作用；
S3--S4检验质量权重；S5--S3检验感知退化影响。
先检查分类别与已见／未见位置的趋势，再结合阶段、失败类型和描述量误差解释收益。
若低质量感知削弱收益，报告其具体影响环节。读出／重构结果用于解释信息保留，不直接估计理论覆盖常数。}
\FloatBarrier

\subsection{History-Conditioned Residual Feedback}
\label{sec:feedbackexperiments}
本节固定同一个RGB＋语义＋QToken基础checkpoint，
检验语义任务上下文与历史能否减少执行错误，并帮助策略从自然偏差中恢复。
表~\ref{tab:feedbackablation}采用MLP／GRU与是否读取语义＋QToken信息的交叉设计；
语义特征和QToken特征在本消融中作为组合加入，不分别归因。

\subsubsection{Matched Inputs and Residual Models}
全部残差组读取最新RGB、本体状态、当前基础计划及目标槽位信息。
不读取语义组合的组也不读取融合语义或QToken后的Transformer特征；
基础动作仍来自共同的语义增强策略。
语义组合组额外读取语义特征与缓存的QToken特征，
分别记录其对应的观测时刻，不将缓存查询视为最新观测。
\draftnote{确认残差语义特征来自本次分割还是缓存；
固定这项实现后再训练所有组合组，报告额外分割带来的实际时延。}

MLP使用当前记录，GRU使用包含相同当前记录的历史窗口。
各组分别训练残差头，匹配目标动作、更新间隔、输出长度、限幅与训练划分；
历史携带的速度或上一指令等信息在各组输入定义中明确。
主要组复核容量接近的当前信息模型，避免仅将参数增加解释为历史收益。

\begin{table*}[t]
\caption{Residual ablations and a correction-method comparison with one frozen semantic-query base policy. Error incidence and recovery are measured from naturally occurring execution errors.}
\label{tab:feedbackablation}
\centering\small
\begin{tabular}{clcccccc}
\toprule
ID & 残差模型 & 语义＋QToken & 历史 & 成功率 & 错误发生率 & 自然恢复率 & 动作MSE\\
\midrule
F0 & 无残差 & 不适用 & 不适用 & -- & -- & -- & --\\
F1 & MLP & 否 & 否 & -- & -- & -- & --\\
F2 & MLP & 是 & 否 & -- & -- & -- & --\\
F3 & GRU & 否 & 是 & -- & -- & -- & --\\
F4 & GRU & 是 & 是 & -- & -- & -- & --\\
\midrule
对照 & A2C2适配 & 匹配 & 原方法 & -- & -- & -- & --\\
\bottomrule
\end{tabular}
\end{table*}

F0与S3使用同一基础策略，F4为完整方法。
另在相同冻结基础策略上接入A2C2适配模块，作为外部修正方法对照，
匹配可用输入与训练预算，并记录其原有逐步修正方式与本文执行方式的实际时延差异。
另以RGB ACT＋RGB历史残差作为连接对照，与S0、S3和F4构成表示增强与反馈的整体比较。
为比较执行方式，在同一基础模型上保留更频繁重规划与temporal ensemble对照，
并记录实际算力和更新频率，而不默认其与残差具有相同计算成本。

\subsubsection{Error Reduction and Natural Recovery}
各组沿用表示消融的物体与初始位置条件，不施加人工扰动。
除分类别、已见／未见位置成功率外，比较阶段完成率和主要失败类型，
并从完整执行视频中记录布料缠阻、无效接触、目标输送偏差和执行停滞四类事件。
这些事件可出现在成功或失败试验中；局部错误后的恢复与最终任务成功分别记录。

错误发生率以出现至少一次错误的试验数除以全部试验数；
自然恢复率以预定窗口内恢复的事件数除以全部已识别事件数，
并报告出错试验的最终成功率和恢复耗时。
无错误事件时恢复率记为不适用，不记为零。
表~\ref{tab:recoveryruntime}按错误类型统计上述结果。
事件起止、恢复条件和重复计数规则见附录~\ref{app:naturalrecovery}。
不同策略经历的错误状态可能不同，因此自然恢复统计与相同基础策略上的总体成功率、错误发生率结合解释。

\begin{table*}[t]
\caption{Natural-error and recovery diagnostics. Instantiate each method separately. Final success is measured over trials containing the specified error; recovery time is conditional on event recovery.}
\label{tab:recoveryruntime}
\centering\small
\begin{tabular}{llcccc}
\toprule
方法 & 错误类型 & 出错试验／总数 & 恢复／事件数 & 出错试验最终成功率 & 恢复时间(s)\\
\midrule
待填 & 布料缠阻 & -- & -- & -- & --\\
待填 & 无效接触 & -- & -- & -- & --\\
待填 & 输送偏差 & -- & -- & -- & --\\
待填 & 执行停滞 & -- & -- & -- & --\\
待填 & 去重合计 & -- & -- & -- & --\\
\bottomrule
\end{tabular}
\end{table*}

\subsubsection{Correction Accuracy and Runtime}
离线纠偏评价使用相同回放样本、专家动作目标、执行槽位和动作尺度，
报告基础动作与修正动作的MSE及式~\eqref{eq:residualgain}中的对齐收益和修正能量。
表~\ref{tab:feedbackruntime}记录端到端延迟、观测年龄、实际更新频率和过期修正比例。
候选修正与实际执行指令分别记录；控制组若保留反馈计算但屏蔽输出，
可以进一步检查计算调度变化是否影响基础策略。

\begin{table*}[t]
\caption{Correction accuracy and runtime. Offline errors use matched evaluation samples; runtime is measured on the deployment device.}
\label{tab:feedbackruntime}
\centering\small
\setlength{\tabcolsep}{4pt}
\begin{tabular}{lcccccc}
\toprule
修正方法 & 对齐收益 & 修正能量 & 延迟p50／p95(ms) & 观测年龄(ms) & 更新频率(Hz) & 过期比例\\
\midrule
F1 & -- & -- & -- & -- & -- & --\\
F2 & -- & -- & -- & -- & -- & --\\
F3 & -- & -- & -- & -- & -- & --\\
F4 & -- & -- & -- & -- & -- & --\\
A2C2适配 & -- & -- & -- & -- & -- & --\\
\bottomrule
\end{tabular}
\end{table*}

\begin{figure*}[t]
\centering
\fbox{\parbox[c][2.5cm][c]{.95\textwidth}{\centering
(a) F0--F4分类别、已见／未见位置成功率及阶段完成率\\[4pt]
(b) 失败类型、出错试验最终成功率及自然恢复结果\\[4pt]
(c) 自然错误--修正--恢复／未恢复视频序列与同步指令轨迹\\[4pt]
标明错误起点、恢复时刻及计划／反馈更新时间；同时展示减少错误与恢复失败案例}}
\caption{Error reduction and natural recovery with residual feedback. Placeholder for task results and synchronized visual and command traces.}
\label{fig:feedbackrecovery}
\end{figure*}

\draftnote{F2--F1与F4--F3检验语义＋QToken组合信息；
F3--F1与F4--F2检验历史。先看同槽动作误差，再看恢复及最终成功，
区分预测改善、修正被及时采用和闭环收益。
同时检查错误发生率：若错误更少，即使恢复事件减少也可能体现反馈收益。
与同基础策略的A2C2比较修正效果和端到端延迟；记录p50／p95、观测年龄与过期比例。
若某组更慢或出现有害修正，结合延迟和修正能量报告，不以更平滑的轨迹代替有效纠偏。
失败机制在Discussion结合独立证据解释。}
\FloatBarrier

\subsection{Cross-Task Applicability}
\label{sec:crosstask}
最后检验相同框架能否支持日常推动以外的操作。
本节在手术模拟任务数据上分别训练各方法，保留语义监督、任务描述量查询和历史残差的整体设计，
评价对象是方法的跨任务适用性，而非同一策略的零样本迁移。

\subsubsection{Surgical-Phantom Task and Adaptation}
该任务使用组织模拟材料遮挡目标，由双臂协同完成牵拉暴露、夹取和向指定区域转移。
相较日常任务，接触方式由推动转为夹持，且暴露动作须与后续抓取协调。
以目标完整进入指定区域并完成释放作为拟定终止条件，
过程评价记录暴露、稳定夹持和放置三个里程碑。
\draftnote{确认最终成功条件是否还要求组织复原，填写时限、持续确认时间、机器人与动作接口、
物体和初始条件、数据规模、独立测试重复以及安全停止规则；图表仅为待完成实验的预留。}

\begin{figure}[t]
\centering
\fbox{\parbox[c][2.8cm][c]{.92\columnwidth}{\centering
(a) 手术模拟平台、器械、相机与材料\\[4pt]
(b) 遮挡--牵拉暴露--稳定夹持--转移放置\\[4pt]
使用同一次试验的关键帧，标明成功终止状态}}
\caption{Surgical-phantom exposure and execution as an additional application. Placeholder for the platform and task sequence.}
\label{fig:surgicaltask}
\end{figure}

各方法使用相同的新任务示范和评价条件，模型选择与测试分开。
需要适配的部分包括机器人动作接口、语义类别以及任务描述量的实例定义；
训练流程与模型设计的共用部分和实际改动分别报告。
\draftnote{列出实际保留与调整的设计、分割标注人工成本和新增训练预算；
若更换描述量、相机布置或控制频率，明确其原因，不将任务相关适配表述为无需修改。}

\subsubsection{Overall Comparison and Key Ablations}
表~\ref{tab:crosstaskresults}比较ACT、DP与完整方法，并保留语义输入、
语义＋QToken无残差、语义＋QToken配合MLP残差三项关键对照。
前两项检查表示收益是否延续，MLP与GRU的比较检查历史是否仍有帮助，
不在此重复日常任务的全部消融。
残差组共用在手术模拟数据上训练的语义查询基础策略。

\begin{table*}[t]
\caption{Cross-task applicability on surgical-phantom execution. Models are trained on this task's demonstrations. Stage entries are cumulative attainment rates; time is conditional on full success.}
\label{tab:crosstaskresults}
\centering\small
\begin{tabular}{p{.30\textwidth}ccccc}
\toprule
方法 & 完整成功率 & 时间(s) & 暴露率 & 夹持率 & 放置率\\
\midrule
ACT & -- & -- & -- & -- & --\\
Diffusion Policy & -- & -- & -- & -- & --\\
ACT＋语义 & -- & -- & -- & -- & --\\
ACT＋语义＋QToken & -- & -- & -- & -- & --\\
ACT＋语义＋QToken＋MLP残差 & -- & -- & -- & -- & --\\
Ours（GRU残差） & -- & -- & -- & -- & --\\
\bottomrule
\end{tabular}
\end{table*}
\draftnote{先比较完整方法与ACT、DP，再检查表示和历史残差的收益是否与日常任务一致。
结合各阶段结果说明提升是否发生在暴露、夹持或放置；若趋势不同，分析任务接触要求与感知条件的差异。
若放置已是全部终止条件，合并表中放置率与完整成功率以避免重复。
结论限于两个应用中实际测试的条件，不表述为未训练策略的跨任务泛化或临床有效性。}
\FloatBarrier
